\id{МРНТИ 52.13.07}{}

\begin{header}
\swa{А.Б.~Жиенбаев, М.А.~Жараспаев, М.Ж.~Балпанова, А.Т.~Карабатырова}{ИССЛЕДОВАНИЕ УСТОЙЧИВОСТИ ГОРНЫХ ВЫРАБОТОК ПРИ ПОДГОТОВКЕ И ТЕХНОЛОГИИ ОТРАБОТКИ ЗАПАСОВ ЗАПАДНОЙ РУДНОЙ ЗОНЫ РУДНИКА АБЫЗ В ЭТАЖЕ 320÷190 М}

\tsp{1}А.Б.~Жиенбаев,
\tsp{2}М.А.~Жараспаев,
\tsp{2}М.Ж.~Балпанова,
\tsp{2}А.Т.~Карабатырова\envelope
\end{header}

\begin{affil}
\tsp{1}ТОО «Корпорация Казахмыс», Жезказган, Казахстан,

\tsp{2}ТОО «Научно-технический центр промышленной безопасности», Караганды, Казахстан

\corrauthor{Корреспондент-автор: akarabatyrova19@gmail.com}
\end{affil}

В данной статье представлено комплексное геомеханическое исследование
устойчивости горных выработок при отработке Западной рудной зоны (ЗРЗ)
рудника Абыз в этаже 320÷190 м. Экономическая целесообразность подземной
геотехнологии на данном объекте напрямую зависит от эффективности
управления состоянием горного массива, особенно в зонах сопряжений и при
повторной разработке запасов. В работе проанализированы
физико-механические свойства пород, которые были разделены на три
основных домена: от крепких сплошных руд до неустойчивых метасоматитов.

Методология исследования базируется на натурных замерах параметров
трещиноватости (742 замера), оценке индекса качества массива RQD и
определении природного поля напряжений методом гидроразрыва скважин. С
использованием программного обеспечения Dips, Unwedge и RocLab проведено
математическое моделирование напряженно-деформированного состояния и
выполнена оценка прочности массива по критерию Хука-Брауна.
Исследованием установлено, что ведущим механизмом потери устойчивости
является кинематическая неустойчивость, проявляющаяся в вывалах
структурных блоков, что значительно усиливается воздействием высоких
горизонтальных тектонических напряжений, превышающих вертикальные в 1,6
раза. Полученные результаты позволяют с высокой точностью прогнозировать
зоны потенциального разрушения и обосновывают необходимость
дифференцированного подхода к оценке устойчивости штреков и ортов в
зависимости от их ориентации относительно главных систем трещин.

{\bfseries Ключевые слова:} устойчивость горных выработок, рудник Абыз,
геомеханический мониторинг, трещиноватость массива, показатель RQD,
критерий Хука-Брауна, кинематическая неустойчивость,
напряженно-деформированное состояние.

\begin{header}
АБЫЗ КЕНІШІНІҢ БАТЫС КЕН АЙМАҒЫНЫҢ 320--190 М ДЕҢГЕЙІНДЕГІ ҚОРЛАРДЫ ДАЙЫНДАУ ЖӘНЕ ӨНДІРУ ТЕХНОЛОГИЯСЫ КЕЗІНДЕГІ ТАУ-КЕН ҚАЗБАЛАРЫНЫҢ ОРНЫҚТЫЛЫҒЫН ЗЕРТТЕУ

\tsp{1}А.Б.~Жиенбаев,
\tsp{2}М.А.~Жараспаев,
\tsp{2}М.Ж.~Балпанова,
\tsp{2}А.Т.~Карабатырова\envelope
\end{header}

\begin{affil}
\tsp{1}«Қазақмыс Корпорациясы» ЖШС, Жезқазған, Қазақстан,

\tsp{2}«Научно-технический центр промышленной безопасности» ЖШС, Қарағанды, Қазақстан,

e-mail: akarabatyrova19@gmail.com
\end{affil}

Бұл мақалада «Абыз» кенішінің Батыс кен аймағын (БКА) 320÷190 м
қабатында игеру кезіндегі дайындық тау-кен қазбаларының орнықтылығына
жүргізілген кешенді геомеханикалық зерттеу нәтижелері ұсынылған. Аталған
объектідегі жерасты геотехнологиясын іске асырудың экономикалық
тиімділігі дайындық қазбаларының айналасындағы тау жынысы сілемінің
күйін басқару тиімділігіне, әсіресе түйісу аймақтарында және қорларды
қайталап өндіру кезінде тікелей байланысты. Жұмыста тау жыныстарының
физикалық-механикалық қасиеттері талданды, олар берік тұтас рудалардан
бастап орнықсыз метасоматиттерге дейінгі үш негізгі геомеханикалық
доменге бөлінді.

Зерттеу әдістемесі жарықшақтық параметрлерін натурлық өлшеуге (742
өлшем), RQD сілемінің сапа индексін бағалауға және ұңғымаларды гидрожару
әдісімен табиғи кернеулер өрісін анықтауға негізделген. Dips, Unwedge
және RocLab бағдарламалық жасақтамасын қолдана отырып,
кернеулі-деформациялық күйге математикалық модельдеу жүргізілді және
Хук-Браун критерийі бойынша сілемнің беріктігі бағаланды. Зерттеу
барысында орнықтылықты жоғалтудың негізгі механизмі құрылымдық
блоктардың құлауымен көрінетін кинематикалық орнықсыздық екені
анықталды, бұл тік кернеулерден 1,6 есе асатын жоғары көлденең
тектоникалық кернеулердің әсерінен айтарлықтай күшейеді. Алынған
нәтижелер ықтимал бұзылу аймақтарын жоғары дәлдікпен болжауға мүмкіндік
береді және негізгі жарықшақтар жүйесіне қатысты бағытына байланысты
штректер мен орттардың орнықтылығын бағалауға сараланған тәсілдің
қажеттілігін негіздейді.

{\bfseries Түйін сөздер:} Тау-кен қазбаларының орнықтылығы, Абыз кеніші,
геомеханикалық мониторинг, сілемнің жарықшақтығы, RQD көрсеткіші,
Хук-Браун критерийі, кинематикалық орнықсыздық, кернеулі-деформациялық
күй.

\begin{header}
STABILITY RESEARCH OF MINE WORKINGS DURING THE DEVELOPMENT AND MINING TECHNOLOGY OF RESERVES OF THE WESTERN ORE ZONE OF THE ABYZ MINE AT THE 320--190 M LEVEL

\tsp{1}A.B.~Zhienbaev,
\tsp{2}M.A.~Zharaspaev,
\tsp{2}M.Zh.~Balpanova,
\tsp{2}A.T.~Karabatyrova\envelope
\end{header}

\begin{affil}
\tsp{1}"Kazakhmys Corporation" LLC, Zhezkazgan, Kazakhstan,

\tsp{2}"Scientific and Technical Center for Industrial Safety" LLC, Karaganda, Kazakhstan,

e-mail: akarabatyrova19@gmail.com
\end{affil}

This article presents the results of a comprehensive geomechanical study
on the stability of development workings during the exploitation of the
Western Ore Zone of the Abyz mine at depths of 190--320 m. The relevance
of the work is driven by the need for scientific justification of the
reasons for the violation of design contours of workings under
conditions of a complex natural stress field and intensive rock mass
jointing. In the course of the work, a detailed analysis of the physical
and mechanical characteristics of rocks was carried out, which were
classified into three main geomechanical domains: from strong massive
ores to weak metasomatites.

The research methodology is based on in-situ measurements of fracturing
parameters (742 measure\-ments), assessment of the Rock Quality
Designation (RQD) index, and determination of the natural stress field
using the hydraulic fracturing method. Using Dips, Unwedge, and RocLab
software, mathematical modeling of the stress-strain state was
performed, and the strength of the rock mass was assessed according to
the Hoek-Brown criterion. The study established that the leading
mechanism of stability loss is kinematic instability, manifested in the
fallouts of structural blocks, which is significantly enhanced by the
impact of high horizontal tectonic stresses exceeding vertical ones by
1.6 times. The results obtained allow for high-accuracy prediction of
potential failure zones and justify the need for a differentiated
approach to assessing the stability of drifts and crosscuts depending on
their orientation relative to the main joint systems.

{\bfseries Keywords:} stability of mine workings, Abyz mine, geomechanical
monitoring, rock mass jointing, RQD index, Hoek-Brown criterion,
kinematic instability, stress-strain state.

\begin{multicols}{2}
{\bfseries Введение.} Актуальность исследования устойчивости горных
выработок на руднике Абыз обусловлена необходимостью обеспечения
промышленной безопасности и экономической эффективности подземных горных
работ при повторной и углублённой отработке запасов в сложных
горно-геологических условиях. В пределах Западной рудной зоны (ЗРЗ)
формируются неблагоприятные сочетания интенсивной трещиноватости,
выраженной анизотропии природного поля напряжений и резкого контраста
физико-механических свойств пород, что приводит к систематическим
нарушениям проектных контуров выработок и значительным потерям
подготовительных горных выработок.

Анализ опыта эксплуатации подземных выработок на руднике Абыз
показывает, что традиционные подходы к оценке устойчивости, основанные
преимущественно на прочностных характеристиках пород и эмпирических
классификациях массива, не в полной мере отражают реальные механизмы
разрушения в условиях развитой структурной неоднородности и тектонически
обусловленных горизонтальных напряжений. В особенности это проявляется в
зонах сопряжений выработок и при проходке в направлениях, неблагоприятно
ориентированных относительно основных систем трещин.

Объектом исследования является горный массив Западной рудной зоны
рудника Абыз в интервале глубин этажа 320÷190 м, характеризующийся
наличием сближенных рудных тел, метасоматически изменённых пород и
сложной тектонической нарушенности.

Предметом исследования являются процессы деформирования и разрушения
горных выработок под воздействием структурных факторов трещиноватости и
анизотропного природного поля напряжений.

Целью работы является обоснование параметров устойчивости горных
выработок на основе комплексного анализа физико-механических свойств
пород, характеристик трещиноватости массива и параметров природного
напряжённого состояния.

Для достижения поставленной цели в работе решаются следующие задачи:

- геомеханическое районирование массива с выделением доменов,
различающихся по прочностным и деформационным характеристикам;

- картирование и статистический анализ систем трещин с оценкой их
влияния на кинематическую устойчивость выработок;

- определение показателя качества массива RQD и геологического индекса
прочности GSI;

- анализ природного поля напряжений и его анизотропии на основе данных
натурных измерений;

- математическое моделирование механизмов неустойчивости горных
выработок.

Гипотеза исследования заключается в том, что доминирующим механизмом
потери устойчивости горных выработок в пределах Западной рудной зоны
является кинематическая неустойчивость, обусловленная перемещением
структурных блоков по системам трещин, усиленная действием высоких
горизонтальных тектонических напряжений, превышающих вертикальные.

Научная значимость работы состоит в формировании уточнённой
геомеханической модели участка недр, учитывающей совместное влияние
структурной неоднородности массива и анизотропного напряжённого
состояния, что позволяет повысить достоверность прогнозирования зон
потенциальных разрушений и обосновать дифференцированный подход к оценке
устойчивости выработок различной ориентации.

{\bfseries Материалы и методы.} \emph{Характеристика объекта исследования.}
Объектом исследования является горный массив Западной рудной зоны (ЗРЗ)
месторождения Абыз в интервале глубин этажа 320÷190 м. Данный участок
характеризуется сложными горно-геологическими условиями, обусловленными
наличием сближенных рудных тел и интенсивной тектоникой.
Золото-колчеданно-полиметаллическое оруденение локализовано в пределах
субмеридиональной тектонической зоны метасоматитов серицит-кварцевого и
хлорит-кварц-серицитового состава, развитых по вулканогенно-осадочным
породам нижнего девонах {[}1{]}.

Важной особенностью объекта является морфология рудной зоны: мощность
зоны метасоматитов достигает 100--200 м, при этом линзы сплошных руд
массивной текстуры, как правило, окаймляются вкрапленными рудами,
развитыми преимущественно в лежачем боку. Общее простирание толщи пород
и тектонических разломов - субмеридиональное, падение - западное,
изменяющееся от крутого (70°) до пологого.

Рудник Абыз отрабатывает две рудные зоны: Восточная - с 2005 г. сначала
карьером, с 2015 г. - шахтой; Западная - с 2021 г. подземным способом
системой поэтажного обрушения (рис.1).
\end{multicols}

{\bfseries Восточная р.з.}

{\bfseries Западная р.з.}

\emph{В}

\emph{З}

{\bfseries Рис.1- Рудные зоны месторождения Абыз в разрезе З -- В}.

\emph{Этаж 320÷190 м выделен красным пунктиром}

\emph{Методология и этапы исследования.} Научная методология работы
базируется на комплексном подходе, включающем натурные измерения,
лабораторные испытания и математическое моделирование.

Определение ключевых механизмов потери устойчивости выработок в условиях
анизотропного поля напряжений и дискретной трещиноватости массива для
прогнозирования зон обрушений.

Основной формой потери устойчивости в исследуемом интервале глубин
является сочетание кинематического перемещения структурных блоков по
трещинам с воздействием высоких горизонтальных тектонических напряжений,
превышающих вертикальные нагрузки.

Этапы исследования:

- геомеханическое районирование: Классификация породного массива на
домены на основе физико-механических свойств;

- cтруктурный анализ: Картирование систем трещин методом линейной съемки
и статистическая обработка данных;

- оценка качества массива: Определение рейтинга RQD (Rock Quality
Designation) по керну и обнажениям;

- анализ напряженного состояния: Измерение природного поля напряжений
методом гидроразрыва и расчет локальных концентраций напряжений вокруг
выработок;

- математическое моделирование: Расчет параметров кинематической
устойчивости (клиновидных вывалов) и оценка прочности массива по
критерию Хука-Брауна;

- описание методов и инструментария.

Для получения достоверных данных использовались следующие методы и
программные комплексы:

- натурные измерения трещиноватости: Проведена линейная съемка методом
массовых замеров. Всего обработано 742 замера элементов залегания трещин
и тектонических разломов на подэтажах 590÷340 м;

- оценка RQD: Выполнена двумя независимыми методами: линейной съемкой
частоты трещин (FF - Fracture Frequency) по стенкам выработок (98
интервалов общей длиной 500 м) и геомеханическим описанием керна 8
разведочных скважин (69 интервалов). Расчет RQD по стенкам выработок
производился по формуле Приста-Хадсона, учитывающей экспоненциальный
закон распределения расстояний между трещинами {[}2{]};

- определение напряжений: Использован метод гидроразрыва пласта
(hydraulic fracturing) в скважинах. Замеры производились на двух
станциях (глубины 408 м и 275 м) для выявления градиента напряжений с
глубиной.

Математическое моделирование и ПО.

- программа Dips использовалась для статистической обработки
ориентировок трещин и выделения главных систем (стереографические
проекции);

- программа RocLab применялась для расчета прочности трещиноватого
массива по критерию Хука-Брауна {[}3{]} и определения геологического
индекса прочности GSI;

- программа Unwedge использовалась для моделирования кинематической
неустойчивости (построения призм обрушения) в зависимости от ориентации
выработок {[}4{]}.

{\bfseries Обсуждение и результаты.} \emph{Физико-механические свойства и
геомеханическая доменизация.}

Анализ лабораторных испытаний проб руд и пород позволил выделить три
основных геомеханических домена, существенно различающихся по своим
прочностным характеристикам. Это разделение критически важно для
понимания дифференцированного поведения выработок.

Домен 1 (Сплошные руды и неизмененные породы): характеризуется высокой
прочностью. Средняя прочность на одноосное сжатие (UCS) составляет 110
МПа с коэффициентом вариации 12\%. Массив данного домена является
жестким (модуль деформации массива E\_m = 22 ГПа) и, как правило,
устойчивым

Домен 2 (Слабоизмененные породы): включает вулканогенно-осадочные
породы, удаленные от рудной зоны. Средняя прочность UCS снижается до 63
МПа.

Домен 3 (Метасоматиты и вкрапленные руды): представляет собой "рубашку"
вокруг рудных тел, состоящую из серицитовых и хлоритовых метасоматитов.
Данный домен является наименее устойчивым: средняя прочность UCS
составляет всего 37 МПа, при этом коэффициент вариации достигает 30\%,
что указывает на высокую неоднородность.

Сравнительный анализ показывает, что прочность сплошных руд в 6 раз
превышает прочность вмещающих метасоматитов, а жесткость массива --- в
22 раза. Геологически это формирует структуру, где жесткие включения
(рудные тела) заключены в податливую, пластичную матрицу метасоматитов,
склонных к размоканию. Подобная контрастность свойств характерна для
многих месторождений Рудного Алтая и является ключевым фактором
перераспределения напряжений.

\emph{Структурный анализ трещиноватости массива.}

Трещиноватость является определяющим фактором устойчивости на
исследуемых глубинах. Обработка 742 замеров в программе Dips (рис.2).
позволила выявить четыре основные системы трещин.

Наиболее распространены и ярко выражены крутопадающие трещины системы 1
со средним азимутом падения DipDir = 265° и углом падения Dip = 68° на
запад, согласные с простиранием (Strata Strike) и падением (Strata Dip)
всей толщи пород и рудных зон. Продольно секущая система трещин 2 имеет
обратное падение на восток со средним азимутом 85° и углом падения 69°.
Выделены система 3 не часто встречающихся пологих трещин со средним
углом падения 12° и азимутом 266°. Скорее всего, система 3 пологих
трещин связана с пологими разломами, по которым произошли сдвиги и ЗРЗ
«отъехала» на запад от ВРЗ. Система 4 поперечно секущих трещин имеет
средний азимут падения 85° и угол падения 69°.

{\bfseries Strata Strike}

{\bfseries Strata Dip}

{\bfseries Рис.2-Полярная диаграмма трещиноватости массива на руднике Абыз.
Средние углы падения (Dip) и азимуты падения (Dip Direction) систем
трещин (Mean Set Planes) в таблице выделены синим}

Длины трещин и расстояния между ними имеют обратно экспоненциальный
закон распределения. По 393 замерам средняя длина трещин 4,0 м; среднее
расстояние между трещинами (по 346 замерам) - 1,2 м.

Важной характеристикой для оценки сопротивления сдвигу является
состояние поверхности трещин.61\% трещин имеют шероховатые поверхности
(Jr = 4), однако 56\% из них заполнены размягчающимся глинистым
материалом (Ja = 4), что существенно снижает сцепление по контактам
блоков. Средневзвешенный показатель шероховатости Jr составляет 3,2, а
показатель изменения стенок трещин Ja - 2,62. Эти параметры были
использованы для расчета индекса GSI по формуле {[}5{]}:

\[GSI = \frac{52\frac{J_{r}}{J_{a}}}{(1 + \frac{J_{r}}{J_{a}})} + \frac{RQD}{2}\]

\emph{Оценка качества массива (RQD) и его вариативность.}

Исследования выявили высокую неоднородность массива по показателю RQD.
По данным линейной съемки в выработках распределение RQD близко к
нормальному со средним значением 54\%. Однако данные кернового бурения
показывают более широкий разброс значений - от 18\% (в зонах
тектонических нарушений) до 70\%, со средним значением 45\% и высоким
коэффициентом вариации (35\%).

\fig{g/image98}[]\fig{g/image99}[]

{\bfseries Рис.3 - Распределения RQD по результатам линейной съемки по
стенкам выработок (слева) и по описанию керна (справа)}

Для моделирования были приняты дифференцированные значения:

- для сплошных руд и окварцованных пород: RQD = 54\% (средняя
трещиноватость);

- для хлорит-серицитовых метасоматитов («рубашки»): RQD = 30\% (сильная
трещиноватость).

Это подтверждает, что метасоматическая зона представляет собой не только
менее прочную, но и значительно более раздробленную среду.

\emph{Природное поле напряжений и его анизотропия.}

Одним из ключевых результатов исследования стало определение параметров
тензора естественных напряжений. Установлено, что напряженное состояние
массива является существенно анизотропным и тектоническим по своей
природе.

Максимальное напряжение (\emph{σ\tsb{1}}) ориентировано
горизонтально по простиранию рудных зон и составляет 1,6 H (H- глубина
от поверхности).

Вертикальное напряжение (\emph{σ\tsb{2}}) соответствует
гравитационной нагрузке (\emph{σ\tsb{2}}=H).

Минимальное напряжение (\emph{σ\tsb{3}}) действует
горизонтально вкрест простирания и составляет 0,7 H.

Таким образом, коэффициент бокового распора по простиранию
(λ\tsb{1} = 1,6) более чем в два раза превышает коэффициент
вкрест простирания (λ\tsb{3} = 0,7). Эта выраженная
анизотропия предопределяет различный характер разрушения выработок
разной ориентации.

\fig{g/image100}[]

{\bfseries Рис.4 - Природное напряженное состояние массива на руднике Абыз
по данным натурных измерений методом гидроразрыва скважин}

\emph{Анализ механизмов неустойчивости.} Кинематическая неустойчивость
(Вывалы блоков)

Моделирование в программном комплексе Unwedge (рис.6, рис.7),
базирующееся на пересечении выявленных систем трещин (1, 2 и 4),
показало, что доминирующим механизмом разрушения является гравитационное
обрушение структурных блоков (клиновидные вывалы рис.5).

{\bfseries Запад}

{\bfseries Рис.5 - Клиновидный вывал в вент. штреке гор.490 м}

\fig{g/image103}[]

{\bfseries Рис.6 - Геометрия возможных вывалов в штреках (слева) и ортах
(справа) по трещинам (см. рис.2): 1 - наиболее развитой системы,
согласной с залеганием рудных тел и всей толщи пород; 2 - продольно
секущей системы с обратным падением; 4 - поперечно секущей системы}

Результаты расчетов вывалов в выработках разных направлений приведены на
рис.7. На этих графиках по горизонтальной оси используется угол встречи
\emph{∆} между азимутом простирания толщи пород и азимутом проходки
выработки. При проходке выработки с отклонением от простирания в висячий
бок угол \emph{∆} имеет знак (+); в лежачий бок - знак (-).

{\bfseries Рис.7 - Объемы вывалов (слева - по месту образования; справа -
суммарные) в выработках разных направлений: 1 - в штреках по простиранию
толщи пород (∆ = 0°); 2 -- в ортах вкрест простирания, по падению толщи
пород (∆ =90°)}

\emph{{\bfseries 1}}

\emph{{\bfseries 2}}

\emph{{\bfseries 1}}

\emph{{\bfseries 2}}

Результаты анализа ориентации выработок выявили четкие закономерности:

- штреки (по простиранию): при совпадении оси выработки с простиранием
рудных тел (угол встречи \emph{∆} ≤ 0±15°) наблюдается максимальный риск
обрушения крупных блоков из кровли. Это обусловлено подсечкой
крутопадающих трещин системы 1 и 2 вдоль длинной оси выработки;

- орты (вкрест простирания): при проходке вкрест простирания риск
вывалов из кровли минимизируется, однако возникает вероятность сдвижения
блоков с бортов выработки;

- оптимальная ориентация: наименьшие объемы расчетных вывалов
наблюдаются при углах встречи оси выработки с простиранием структур в
диапазоне 20°÷70°.

Важно отметить, что моделирование без масштабирования трещин
предсказывает возможность формирования особо крупных вывалов при
неблагоприятном пересечении длинных магистральных трещин, что
подтверждается практикой (наблюдались вывалы высотой до 4--5 м).

\emph{Стресс-индуцированная неустойчивость (Раздавливание)}

Для оценки риска разрушения под действием сжимающих напряжений был
выполнен расчет концентрации напряжений на контуре выработок с
использованием решения Кирша {[}6{]}. Максимальные напряжения на контуре
\emph{maxσ} определяются по формуле:

\emph{maxσ = γН(3 - λ\tsb{3}) = 2,3γН.}

где λ - коэффициенты концентрации.

Расчеты для горизонта 320--190 м показали критическую разницу для
выработок разной ориентации:

- в штреках (вдоль \emph{σ\tsb{1}}) максимальные напряжения
концентрируются в бортах и составляют \emph{maxσ \textasciitilde{}}2,3H
(до 38 МПа на глубине 610 м);

- в ортах (вкрест \emph{σ1)} максимальные напряжения возникают в кровле
и достигают экстремальных значений \emph{maxσ \textasciitilde{}} 3,8 H
(до 63 МПа) {[}7{]}.

Сравнивая действующие напряжения с прочностью массива
(\emph{σ\tsb{m}}), можно сделать вывод, что в кровле ортов
напряжения на 66\% выше, чем в бортах штреков.

При прочности массива метасоматитов \emph{σ\tsb{m}}
\textasciitilde6 МПа (табл.2), условие \emph{maxσ}\textasciitilde{}
\emph{σ\tsb{m}} выполняется с большим запасом, что
прогнозирует процессы раздавливания и пластического течения пород в
кровле поперечных выработок, что согласуется с современными
представлениями об инициации разрушения хрупких горных пород в условиях
концентрации напряжений {[}8{]}.

{\bfseries Таблица 1 - Свойства разных типов трещиноватых массивов
(доменов)}

%% \begin{longtable}[]{@{}
%%   >{\raggedright\arraybackslash}p{(\linewidth - 18\tabcolsep) * \real{0.0500}}
%%   >{\raggedright\arraybackslash}p{(\linewidth - 18\tabcolsep) * \real{0.2641}}
%%   >{\centering\arraybackslash}p{(\linewidth - 18\tabcolsep) * \real{0.0927}}
%%   >{\centering\arraybackslash}p{(\linewidth - 18\tabcolsep) * \real{0.0824}}
%%   >{\centering\arraybackslash}p{(\linewidth - 18\tabcolsep) * \real{0.0905}}
%%   >{\centering\arraybackslash}p{(\linewidth - 18\tabcolsep) * \real{0.0954}}
%%   >{\centering\arraybackslash}p{(\linewidth - 18\tabcolsep) * \real{0.0873}}
%%   >{\centering\arraybackslash}p{(\linewidth - 18\tabcolsep) * \real{0.0783}}
%%   >{\centering\arraybackslash}p{(\linewidth - 18\tabcolsep) * \real{0.0796}}
%%   >{\centering\arraybackslash}p{(\linewidth - 18\tabcolsep) * \real{0.0797}}@{}}
%% \toprule\noalign{}
%% \multicolumn{2}{@{}>{\centering\arraybackslash}p{(\linewidth - 18\tabcolsep) * \real{0.3141} + 2\tabcolsep}}{%
%% \begin{minipage}[b]{\linewidth}\centering
%% {\bfseries Домены}
%% \end{minipage}} & \begin{minipage}[b]{\linewidth}\centering
%% {\bfseries UCS}
%% \end{minipage} & \begin{minipage}[b]{\linewidth}\centering
%% {\bfseries J\tsb{r}}
%% \end{minipage} & \begin{minipage}[b]{\linewidth}\centering
%% {\bfseries J\tsb{a}}
%% \end{minipage} & \begin{minipage}[b]{\linewidth}\centering
%% {\bfseries RQD}
%% \end{minipage} & \begin{minipage}[b]{\linewidth}\centering
%% {\bfseries GSI}
%% \end{minipage} & \begin{minipage}[b]{\linewidth}\centering
%% m
%% \end{minipage} & \begin{minipage}[b]{\linewidth}\centering
%% {\bfseries σ\tsb{m}}
%% \end{minipage} & \begin{minipage}[b]{\linewidth}\centering
%% {\bfseries E\tsb{m}}
%% \end{minipage} \\
%% \midrule\noalign{}
%% \endhead
%% \bottomrule\noalign{}
%% \endlastfoot
%% 1 & руды сплошные, туфы, песчаники и другие неизмененные породы & 110 &
%% 4,0 & 0,75 & 54 & 70 & 17 & 38 & 22 \\
%% 2 & слабоизмененные вулканогенно-осадочные породы разных типов & 63 &
%% 3,2 & 2,6 & 45 & 50 & 17 & 14 & 6 \\
%% 3 & слабые метасоматиты серицитовые, хлоритовые, вкрапленные руды & 37 &
%% 2,0 & 4,0 & 30 & 32 & 19 & 6 & 1 \\
%% \end{longtable}

\emph{UCS - прочность руды / пород при сжатии в образце, МПа; m --
показатель Хука, характеризующий тип породы;}

\emph{σ\tsb{m} - прочность трещиноватого массива при сжатии,
МПа; E\tsb{m} -- модуль деформации трещиноватого массива,
ГПа.}

\emph{Валидация модели фактическими данными.} Анализ фактического
состояния горных выработок полностью подтверждает теоретические и
модельные построения.

За период эксплуатации подземного рудника из-за обрушений потеряно 1693
погонных метра выработок, что составляет значительный экономический
ущерб. Доминирование записей "Обрушение" в актах обследования (табл.2)
коррелирует с прогнозом высокой кинематической неустойчивости.

\tcap{Таблица 2 - Бросовые выработки на руднике Абыз на 01.03.2024 г.}

%% \begin{longtable}[]{@{}
%%   >{\centering\arraybackslash}p{(\linewidth - 4\tabcolsep) * \real{0.2186}}
%%   >{\raggedright\arraybackslash}p{(\linewidth - 4\tabcolsep) * \real{0.5783}}
%%   >{\raggedright\arraybackslash}p{(\linewidth - 4\tabcolsep) * \real{0.2030}}@{}}
%% \toprule\noalign{}
%% \begin{minipage}[b]{\linewidth}\centering
%% {\bfseries Горизонт}
%% \end{minipage} & \begin{minipage}[b]{\linewidth}\raggedright
%% {\bfseries Выработки (в скобках длина в п.м.)}
%% \end{minipage} & \begin{minipage}[b]{\linewidth}\raggedright
%% {\bfseries Итого (п.м.)}
%% \end{minipage} \\
%% \midrule\noalign{}
%% \endhead
%% \bottomrule\noalign{}
%% \endlastfoot
%% {\bfseries I. Обрушение} & & 1693 \\
%% 657 & БПШ-1 (35), БПШ-4 бис (15), БДШ-1 (70), БПШ-5 бис (45) & 165 \\
%% 640 & БПШ-10 (60), БПШ-8 (55), БПШ-5 (40), БПШ-4 (35), БПШ-1 (75), БПШ-2
%% (30), БПШ-3 (25) & 320 \\
%% Склад ВМ & Склад ВМ (110) & 110 \\
%% 623 & ОРТ-3 (55), РШ-2 (15), Заезд-1 (15) & 85 \\
%% 607 & БПШ-2 (125) & 125 \\
%% 540 & РШ-1 (44), РШ-2 (33), РШ-3 (23) & 100 \\
%% 457 & БПШ-6 (48), ОРТ-3 (20) & 68 \\
%% 440 & БПШ-1 (25) & 25 \\
%% 423 & Заезд на БПШ-4 (20), БПШ-8 (30), ВОШ (25), БПШ-3 (66) & 141 \\
%% 406 & БПШ-7 (38), БПШ-5 (30), БПШ-5 бис (48), БПШ-4 бис (48), БПШ-6 бис
%% (4) & 168 \\
%% 390 & БПШ-2 (47), БПШ-4 (31), БПШ-1 (76), БПШ-3 (29), БПШ 1бис (15),
%% БПШ-1 (40) & 238 \\
%% 373 & БПШ-1 (61), Подходной штрек (15), БПШ-1 (16) & 92 \\
%% 356 & БПШ-1 (9) & 9 \\
%% 340 & БПШ-1 (19), ОРТ-5 (19), ОРТ-4 (9) & 47 \\
%% {\bfseries II.Неподтверждение руды} & & 848 \\
%% 623 & БПШ-11 бис (90) & 90 \\
%% 590 & БПШ-7 (35), Заезд на ВВ (35), РШ-2 (31), РШ-1 (16), БПШ-5 (211) &
%% 328 \\
%% 573 & БПШ-1 (11), БПШ-2 (19), БПШ-3 (66), Рудный штрек (78) & 174 \\
%% 557 & Заезд на п/э 557 м. (22) & 22 \\
%% 540 & Заезд на п/э 540 м. (20) & 20 \\
%% 523 & Заезд на п/э 523 м. (25) & 25 \\
%% 490 & БПШ-2 (25) & 25 \\
%% 473 & БПШ-7 (58) & 58 \\
%% 423 & БПШ-1 (61) & 61 \\
%% 406 & РШ-4 (45) & 45 \\
%% Итого по руднику & & 2541 \\
%% \end{longtable}

\fig{g/image105}[]В
актах обследований и при визуальном осмотре фиксируется крупноблочный
состав обрушенной массы, что является прямым признаком кинематического
механизма (структурной неустойчивости), а не дробления давлением.

{\bfseries Рис.8 - Вывалы горной массы в БПО п/э 340 м ЗРЗ}

Наблюдаемые куполообразные вывалы высотой 1,0--1,5 м (в редких случаях
до 5 м) соответствуют геометрии призм обрушения, построенных в Unwedge
(рис.9).

\fig{g/image107}[Рис. 9 - Клиновидные вывалы с бортов и кровли БПШ-1 на п/э 390 м в сплошной
медной колчеданной руде]

Инструментальные замеры показали, что средний перебор сечения составляет
3,3 м² (23\% от проектного), что свидетельствует о самопроизвольном
обрушении приконтурного массива сразу после взрывных работ до момента
установки крепи {[}9{]}. Это указывает на мгновенную реализацию упругой
энергии и раскрытие трещин в условиях высоких напряжений. Результаты 48
промеров приведены на рис.10.

\fig{g/image108}[Рис.10 - Распределение площадей переборов сечений выработок]

\begin{multicols}{2}
Установленные параметры переборов и объемы обрушений являются основанием
для корректировки существующих паспортов крепления в соответствии с
действующим Регламентом применения горных мер обеспечения устойчивости
{[}10{]}.

{\bfseries Выводы.} Установлено, что в пределах Западной рудной зоны
рудника Абыз в интервале глубин этажа 320--190 м определяющим механизмом
потери устойчивости горных выработок является кинематическая
неустойчивость, реализующаяся в форме гравитационных вывалов структурных
блоков, ограниченных системами естественных трещин. Данный механизм
доминирует над разрушением, обусловленным превышением прочности массива
при сжатии.

Показано, что кинематическая неустойчивость существенно усиливается в
зонах метасоматически изменённых пород, формирующих «рубашку» вокруг
рудных тел, для которых характерны низкие значения показателя качества
массива (RQD ≈ 30\%), высокий уровень структурной нарушенности и наличие
размягчающего глинистого заполнителя в трещинах (Ja = 4). В этих
условиях даже при умеренных напряжениях реализуются крупноблочные
вывалы.

Экспериментально подтверждено наличие выраженного анизотропного
тектонического поля напряжений, при котором максимальное горизонтальное
напряжение ориентировано по простиранию рудной зоны и превышает
вертикальное в 1,6 раза. Такая конфигурация напряжённого состояния
формирует принципиально различный характер нагружения выработок в
зависимости от их пространственной ориентации.

Установлено, что для ортов, пройденных вкрест простирания рудных тел, в
кровле формируются экстремальные концентрации сжимающих напряжений (до
3,8H), что создаёт предпосылки для раздавливания и пластического течения
пород даже в доменах с относительно высокой прочностью. В штреках,
ориентированных по простиранию, максимальные напряжения концентрируются
преимущественно в бортах и достигают уровня около 2,3H.

Показана прямая зависимость объёмов потенциальных вывалов от угла
встречи оси выработки с основными системами трещин. Наиболее
неблагоприятным является направление проходки, совпадающее с
простиранием пород и рудных тел (α ≈ 0°), при котором создаются условия
для формирования крупных клиновидных вывалов из кровли. Минимальные
расчётные объёмы неустойчивых блоков характерны для диапазона углов
20°--70°.

Сопоставление результатов математического моделирования с фактическими
данными эксплуатации подтвердило адекватность принятой геомеханической
модели: наблюдаемые куполообразные и клиновидные вывалы, значительные
переборы сечений (в среднем 23\%) и протяжённость потерянных выработок
(1693 м) полностью коррелируют с прогнозируемыми зонами кинематической и
стресс-индуцированной неустойчивости.

Полученные результаты свидетельствуют о том, что применение
унифицированных подходов к обеспечению устойчивости горных выработок без
учёта структурного фактора и анизотропии напряжений является
неэффективным в условиях Западной рудной зоны рудника Абыз. Обоснована
необходимость дифференцированного геомеханического подхода, основанного
на совместном учёте прочностных характеристик пород, параметров
трещиноватости, ориентации выработок и тензора природных напряжений.

Практическая реализация предложенных выводов предполагает корректировку
направлений проходки, уточнение паспортов крепления и внедрение
инструментального мониторинга деформаций, что позволит снизить риск
потери выработок и повысить технико-экономическую эффективность
подземной отработки запасов.
\end{multicols}

\begin{center}
{\bfseries Литература}
\end{center}

\begin{refs}
1. Разработка ресурсосберегающей технологии крепления
горно-подготовительных выработок в геомеханических условиях
месторождения Абыз {[}док. внутреннего пользования{]}/ ТОО «Expert
PRO». - Усть-Каменогорск, 2023. - 408 с.

2. Hudson J.A., Harrison J.P. Engineering Rock Mechanics. An
Introduction to the Principles/ Oxford: Elsevier Science Ltd. -2000.
-456 p. ISBN 0-08-041912-7, 0-08-043864-4.

3. Hoek E. Practical Rock Engineering. -2007. -341p. URL:
\url{https://www.rocscience.com/learning/hoek-s-corner/books}. -Дата
обращения: 14.11.2025.

4. Using the Q-system. Rock mass classification and support design:
NGI Handbook/Oslo: NGI. -2025.
\url{https://www.ngi.no/globalassets/bilder/forskning-og-radgivning/bygg-og-anlegg/q-handbook_2025_english_final.pdf}.
-Дата обращения: 14.11.2025.

5. Рекомендации по параметрам отработки подкарьерных запасов на
руднике Абыз: ГМО № 01-9.4.3.5-9--290 от 04.11.2020. {[}док.
внутреннего пользования{]}/ТОО «Корпорация Казахмыс». - Караганда,
2020. - 32 с.

6. Макаров А. Б. Практическая геомеханика. - М.: Горная книга, 2006.-
391 с. ISBN 5-98672-038-5.

7. Баклашов, И. В. Геомеханика. В 2 т. Т.1. Основы геомеханики:
учебник для вузов. - М.: Изд-во МГГУ, 2004. -- 208 с. ISBN
5-7418-0327-Х, ISBN 5-7418-0325-3(Т.1).

8. Diederichs M. S., Kaiser P. K., Eberhardt E. Damage initiation and
propagation in hard rock during tunnelling and the influence of
near-face stress rotation // International Journal of Rock Mechanics
and Mining Sciences. - 2004.-Vol.41(5). -P.785-812.
\href{https://doi.org/10.1016/j.ijrmms.2004.02.003}{DOI
10.1016/j.ijrmms.2004.02.003}.

9. Мангуш, С. К. Взрывные работы при проведении подземных горных
выработок: учебное пособие / С. К. Мангуш. - 2-е изд., стер. - Москва:
Издательство МГГУ, 2005. - 120 с. ISBN 5-7418-0114-5.

10. Регламент применения горных мер обеспечения устойчивости выработок на
рудниках Корпорации Казахмыс: ГМО № 01--7.10.2.4-9--197 от 27.07.2022 г.
{[}док. внутреннегоего пользования{]}/ТОО «Корпорация Казахмыс». -
Караганда, 2024.
\end{refs}

\begin{center}
{\bfseries References}
\end{center}

\begin{refs}
1. Razrabotka resursosberegajushhej tehnologii kreplenija
gorno-podgotovitel' nyh vyrabotok v geomehanicheskih
uslovijah mestorozhdenija Abyz {[}dok. vnutrennego
pol' zovanija{]}/ TOO «Expert PRO». -
Ust'-Kamenogorsk, 2023. - 408 s. {[}in Russian{]}

2. Hudson J.A., Harrison J.P. Engineering Rock Mechanics. An
Introduction to the Principles/ Oxford: Elsevier Science Ltd. -2000.
-456 p. ISBN 0-08-041912-7, 0-08-043864-4.

3. Hoek E. Practical Rock Engineering. -2007. -341p. URL:
\url{https://www.rocscience.com/learning/hoek-s-corner/books}. - Date
accessed: 14.11.2025.

4. Using the Q-system. Rock mass classification and support design: NGI
Handbook/Oslo: NGI. -2025.
\url{https://www.ngi.no/globalassets/bilder/forskning-og-radgivning/bygg-og-anlegg/q-handbook_2025_english_final.pdf}.
- Date accessed: 14.11.2025.

5. Rekomendacii po parametram otrabotki podkar' ernyh
zapasov na rudnike Abyz: GMO № 01-9.4.3.5-9-290 ot 04.11.2020. {[}dok.
vnutrennego pol' zovanija{]}/TOO «Korporacija Kazahmys. -
Karaganda, 2020. - 32 s. {[}in Russian{]}

6. Makarov A. B. Prakticheskaja geomehanika. - M.: Gornaja kniga, 2006.-
391 s. ISBN 5-98672-038-5. {[}in Russian{]}

7. Baklashov, I. V. Geomehanika. V 2 t. T.1. Osnovy geomehaniki:
uchebnik dlja vuzov. - M.: Izd-vo MGGU, 2004. -- 208 s. ISBN
5-7418-0327-H, ISBN 5-7418-0325-3(T.1). {[}in Russian{]}

8. Diederichs M. S., Kaiser P. K., Eberhardt E. Damage initiation and
propagation in hard rock during tunnelling and the influence of
near-face stress rotation // International Journal of Rock Mechanics and
Mining Sciences. - 2004.-Vol.41(5). -P.785-812.
\href{https://doi.org/10.1016/j.ijrmms.2004.02.003}{DOI
10.1016/j.ijrmms.2004.02.003}.

9. Mangush, S. K. Vzryvnye raboty pri provedenii podzemnyh gornyh
vyrabotok: uchebnoe posobie / S. K. Mangush. - 2-e izd., ster. - Moskva:
Izdatel' stvo MGGU, 2005. - 120 s. ISBN 5-7418-0114-5.
{[}in Russian{]}

10. Reglament primenenija gornyh mer obespechenija ustojchivosti
vyrabotok na rudnikah Korporacii Kazahmys: GMO № 01--7.10.2.4-9--197 ot
27.07.2022 g. {[}dok. vnutrennegoego pol' zovanija{]}/TOO
«Korporacija Kazahmys». - Karaganda, 2024. {[}in Russian{]}
\end{refs}

\begin{info}
\hspace{1em}\emph{{\bfseries Сведения об авторах}}

Жиенбаев А.Б. - PhD, заместитель начальника отдела геомеханики, ТОО
«Корпорация Казахмыс», Жезказган, Казахстан, e-mail:
zhienbaev@list.ru;

Жараспаев М.А. - ведущий специалист, ТОО «Научно-технический центр
промышленной безопасности», Караганда, Казахстан, e-mail:
zharaspaev\_m\_a@mail.ru;

Балпанова М. Ж. - PhD, директор ТОО «Научно-технический центр
промышленной безопасности», Караганда, Казахстан, e-mail:
balpanova86@mail.ru;

Карабатырова А. Т. - ведущий специалист, ТОО «Научно-технический центр
промышленной безопасности», Караганда, ККазахстан,e- mail
akarabatyrova19@gmail.com.

\hspace{1em}\emph{{\bfseries Information about the authors}}

Zhienbaev A.B. - PhD, Deputy Head of the Geomechanical Department,
"Kazakhmys Corporation" LLC, Zhezkazgan, Kazakhstan, e-mail:
zhienbaev@list.ru;

Zharaspaev M.A. - Leading Specialist, "Scientific and Technical Center
for Industrial Safety" LLC,Karaganda, Kazakhstan, e-mail:
zharaspaev\_m\_a@mail.ru;

Balpanova M. Zh. - PhD, Director, "Scientific and Technical Center for
Industrial Safety" LLC, Karaganda, Kazakhstan, e-mail:
balpanova86@mail.ru;

Karabatyrova A.T. - Leading Specialist, "Scientific and Technical Center
for Industrial Safety" LLC, Karaganda, Kazakhstan, e-mail:
aiger\_aiger95@mail.ru.
\end{info}
